\scanpage{322}

\mspage{61}
% archive sheet number: 321

On a alors
\[
f_{+2}^{*}\bigl(\underbrace{a\overset{u}{\text{ --- }}b\overset{v}{\text{ --- }}c}_{c(v)c(u)}\bigr)
=\underset{c(vu)}{(a\xrightarrow{vu}c)}
\]
d'où une \unsure{flèche|flèches} canonique
\[
c(vu)\longrightarrow c(v)c(u)\qquad\text{dans }\underline{\mathrm{Ch}}(a,c)
\]

Cela montre que $c(vu)$ et $c(v)c(u)$ sont \unsure{donc} au moins dans une même composante connexe de $\underline{\mathrm{Ch}}(a,c)$, ou \deleted{\unsure{encore}} (comme on dit aussi) qu'ils sont ``homotopes''. Cela suggère \inserted{\unsure{ici}} de \deleted{\unsure{déf}} modifier la définition des $\operatorname{Hom}_{\mathrm{ch}}(a,b)$ en
\[
\pi(a,b)=\pi_0\bigl(\underline{\mathrm{Ch}}(a,b)\bigr).
\]

Le foncteur de composition des chemins
% slanted
\margin{Ce foncteur est déduit des foncteurs $\underline{\mathrm{Ch}}_\tau(a,b)\times\underline{\mathrm{Ch}}_{\tau'}(b,c)\longrightarrow\underline{\mathrm{Ch}}_{\tau\tau'}(a,c)$ \unsure{evidents|évidents}, cf \unsure{pag|pge} \unsure{7.7|77|47}}
\[
\begin{aligned}
\underline{\mathrm{Ch}}(a,b)\times\underline{\mathrm{Ch}}(b,c)&\;\text{ --- }\;\underline{\mathrm{Ch}}(a,c)\\
(\unsure{a|c},c')&\;\unsure{\mathrel{\mapstochar\relbar}|\longmapsto}\;c'\circ c
\end{aligned}
\]
induit une \deleted{\unsure{``compos}} application sur les $\pi_0$, ce qui donne une loi de composition
\[
\pi(a,b)\times\pi(b,c)\;\text{ --- }\;\pi(a,c)
\]
sur les flèches. Il est immédiat qu'elle est associative (la composition des chemins l'étant), et unitaire. Elle admet comme flèche identique en $x$ soit le \deleted{\unsure{mot}} \deleted{chemin} ``le mot vide en $x$'', lequel est \inserted{(évidemment)} homotope dans $\underline{\mathrm{Ch}}(x,x)$ au \unsure{chemin} ``dégénéré'' $x\xrightarrow{\mathrm{id}}x$
